\section{Modeling Objective}
\label{sec:8-1}

Estimate $P(\text{Churn} = 1 \mid \text{Customer History at Snapshot})$.

\section{Baselines}
\label{sec:8-2}

\begin{center}
\begin{tabular}{@{}p{4.2cm}p{9.9cm}@{}}
    \toprule
    Baseline & Role \\
    \midrule
    Random ranking & Reference for Lift $= 1$ and PR-AUC $=$ churn rate. \\
    Recency rule & Rank customers by recency (longest since last purchase first). A strong, business-realistic rule the models must beat. \\
    Optional: RFM score rule & Rank by a simple combined RFM score. \\
    \bottomrule
\end{tabular}
\end{center}

\section{RFM Model}
\label{sec:8-3}

A model using Recency, Frequency and Monetary features.

\section{Extended Feature Model}
\label{sec:8-4}

Add behavioral and temporal features to test whether richer history improves prediction.

\section{Full Feature Model}
\label{sec:8-5}

Add product and cancellation features where supported.

\section{Algorithms}
\label{sec:8-6}

Compare Logistic Regression, Random Forest and XGBoost.

\section{Hyperparameter Tuning}
\label{sec:8-7}

Tune on the validation period while preserving temporal separation. Validation metrics are computed on the Scoring Population. After selection, \textbf{retrain the chosen configuration on the Training Population re-derived over train + validation} (respecting the embargo before the test period) and evaluate once on the test Scoring Population.

\section{Model Selection}
\label{sec:8-8}

Select the final model using predictive and ranking performance, with particular attention to Top-K targeting metrics at the operating $K$.

\section{Model Interpretability}
\label{sec:8-9}

Feature importance and, where appropriate, global/local explanations (e.g.\ SHAP). Findings are described as associations, not causes.

\section{Training Population Ablation}
\label{sec:8-10}

Compare two training schemes for the selected algorithm and feature set, both evaluated on the \textbf{same} test Scoring Population:

\begin{center}
\begin{tabular}{@{}p{4.5cm}p{9.6cm}@{}}
    \toprule
    Scheme & Training rows \\
    \midrule
    A -- Episode (main) & \texttt{in\_training\_sample} rows, unweighted \\
    B -- All weekly, weighted & All weekly active rows in the training period, weighted by \texttt{sample\_weight} \\
    \bottomrule
\end{tabular}
\end{center}

If A matches B on PR-AUC and weekly Lift@K, the episode rule reduces training data and redundancy without losing predictive quality. If B is clearly better, report it and discuss the trade-off.
